% ============================================================================
% PATCH v5.8.0.02: κ as First-Class Open Problem
% ============================================================================
% Issue: κ bridge is Achilles' heel, blocks absolute predictions
% Fix: Formalize as Grand Challenge Problem, not bridge-in-waiting
% ============================================================================

\section{Grand Challenge Problem: The κ Bridge}
\label{sec:kappa-grand-challenge}

\begin{tcolorbox}[colback=red!5!white,colframe=red!75!black,title=Grand Challenge Problem \#1: Deriving κ]
\label{gcp:kappa}
\textbf{Problem:} Derive the conversion constant $\kappa$ in $\varphi = \kappa \tau_{\mathrm{lat}}$ from first principles of the substrate dynamics, or provide an independent experimental protocol to measure it.

\textbf{Current status:} [Open Problem / Ansatz]

\textbf{Why it matters:} Until resolved, all absolute predictions (e.g., $G$, timing effects) remain contingent.
\end{tcolorbox}

\subsection{What We Know vs. What We Need}

\begin{tabular}{p{0.45\textwidth}p{0.45\textwidth}}
\textbf{Currently Established} & \textbf{Required for Resolution} \\ \hline
\textbf{Dimensional analysis:} $[\kappa] = L^2 T^{-2} (cy/b)^{-1}$ & \textbf{First-principles derivation:} $\kappa = f(\text{substrate parameters})$ \\
\textbf{Formal definition:} $\kappa = \hbar f_U / (m_P c^2)$ & \textbf{Independent measurement:} Protocol for $f_U$ \\
\textbf{Consistency:} Recovers Newtonian limit when matched & \textbf{Prediction:} Use derived $\kappa$ to predict new observable \\
\textbf{Honest labeling:} Clearly marked as ansatz & \textbf{Falsification:} Clear test if derivation fails \\
\end{tabular}

\subsection{Consequences of Unresolved κ}

\begin{enumerate}
\item \textbf{No absolute predictions:} Any claim depending on $\kappa$'s numerical value cannot be labeled [Theorem].

\item \textbf{Circularity risk:} Using $\kappa$ to "derive" $G$ then checking against measured $G$ is circular unless $\kappa$ is fixed independently.

\item \textbf{Propagation of uncertainty:} Downstream claims inherit κ's ansatz status.

\item \textbf{Critical vulnerability:} Critics correctly state "all absolute predictions are contingent."
\end{enumerate}

\subsection{Roadmap for Resolution}

\subsubsection{Path A: Theoretical Derivation}
\begin{enumerate}
\item \textbf{Substrate dynamics:} Derive $f_U$ from Leech lattice vibrational spectrum.
\item \textbf{Information flow:} Relate mapping cycles to Planck-scale computations.
\item \textbf{Emergent units:} Show how substrate computations define $L$ and $T$.
\item \textbf{Prediction:} Compute $\kappa$, then predict $G = (8\pi)^{-1} \kappa^2 / c^4$.
\item \textbf{Test:} Match against measured $G = 6.67430(15) \times 10^{-11} \text{m}^3\text{kg}^{-1}\text{s}^{-2}$.
\end{enumerate}

\subsubsection{Path B: Experimental Measurement}
\begin{enumerate}
\item \textbf{Substrate signatures:} Look for $f_U \sim 10^{43}$ Hz in:
\begin{itemize}
\item Vacuum fluctuation spectra
\item Planck-scale discreteness in collisions
\item Anomalous timing in atomic clocks
\end{itemize}

\item \textbf{Independent prediction:} Use measured $f_U$ to compute $\kappa$, then predict:
\begin{itemize}
\item Variation of $G$ with computational load
\item Quantum gravity corrections to black hole entropy
\item Specific CMB non-Gaussianity from discrete sampling
\end{itemize}

\item \textbf{Falsification:} If $\kappa$ measured this way gives wrong $G$, framework falsified.
\end{enumerate}

\subsubsection{Path C: Reformulation}
\begin{enumerate}
\item \textbf{Dimensionless focus:} Emphasize predictions where $\kappa$ cancels.
\item \textbf{Relative predictions:} Predict ratios, not absolutes.
\item \textbf{Qualitative claims:} Make claims independent of $\kappa$'s value.
\end{enumerate}

\subsection{Interim Protocol for κ-Dependent Work}

Until resolved, follow this protocol for any claim $C$ depending on $\kappa$:

\begin{tcolorbox}[colback=yellow!5!white,colframe=yellow!75!black,title=Protocol: Handling κ-Dependent Claims]
\begin{enumerate}
\item \textbf{Label:} $C$ must be labeled [Ansatz] or [Conjecture], never [Theorem].
\item \textbf{Dependencies:} Must list "κ ansatz" as explicit dependency.
\item \textbf{Track:} Must be in Track B (Calibrated) or Track C (Speculative).
\item \textbf{Falsifiability:} Must provide test independent of κ matching.
\item \textbf{Transparency:} Must state "contingent on resolution of Grand Challenge \#1."
\end{enumerate}
\end{tcolorbox}

\subsection{Relation to Other Open Problems}

\begin{itemize}
\item \textbf{Leech uniqueness:} If substrate not unique, $\kappa$ derivation ambiguous.
\item \textbf{Substrate dynamics:} Need equations to compute $f_U$.
\item \textbf{Emergent gravity:} κ connects informational and geometric descriptions.
\item \textbf{Testability:} Resolution enables absolute experimental tests.
\end{itemize}

\subsection{Success Criteria}

The κ problem is resolved when \textbf{either}:
\begin{enumerate}
\item $\kappa$ is derived from substrate parameters and predicts $G$ within experimental error, \textbf{or}
\item $f_U$ is measured independently and yields correct $G$, \textbf{or}  
\item All framework predictions are reformulated to be κ-independent.
\end{enumerate}

\paragraph{Current Assessment}
As of v5.8.0, κ remains the \textbf{primary obstacle} to the framework making absolute, falsifiable predictions. It is not a minor technical gap but a central conceptual challenge that must be addressed for the framework to achieve full scientific credibility.
